\par\medskip\pdfbookmark[1]{On III and IV.}{nexus-note-III-IV}\centerline{On III and IV.}\smallskip
Let $C$ be the collection of all positive, properly primitive forms of negative determinant $-D$ that are not properly equivalent. Assign to the variables of each form two values relatively prime to one another. The total number of representations of a positive integer $m$ is then $\epsilon\psi(m)$, where $\epsilon$ denotes the number of solutions of the equation $tt+Duu=1$, and $\psi(m)$ the number of roots $n$ of the congruence $nn+D\equiv0\pmod m$ for which the three numbers $m,\ 2n$ and $\frac{nn+D}m$ have no common divisor (Disqq. Arithm., art. 180). For the factor, $\epsilon$$=4$ when $D=1$, and $=2$ in every other case. Further, if $m=p^\pi p'^{\pi'}p''^{\pi''}\ldots$, where $p,p',p''\ldots$ are distinct primes, then $\psi(m)=\psi(p^\pi)\psi(p'^{\pi'})\psi(p''^{\pi''})\ldots$; if $\mathfrak A(m)$ denotes the number of roots $n$ of the congruence $nn+D\equiv0\pmod m$, and one uses the symbol introduced by \textsc{Legendre} and generalized by \textsc{Jacobi}, then $\psi(p^\pi)=\mathfrak A(p^\pi)=1+\left(\frac{-D}p\right)$ when $p$ does not divide $2D$, but otherwise $=\mathfrak A(p^\pi-\frac1p\mathfrak A(p^{\pi+1})$; the number $\mathfrak A(p^\pi)$ is always easy to determine (Disqq. Arithm., art. 104). For what follows, however, it suffices to observe that $\mathfrak A(p^\pi)$ always becomes independent of $\pi$ once $\pi$ exceeds a certain size.

If the variables of the forms in the collection $C$ are assigned all integer values without exception (Disqq. Arithm., art. 181), the total number $(m)$ of representations of the number $m$ becomes $\epsilon f(m)$, where $f(m)=\Sigma\psi\left(\frac m{\mu\mu}\right)$ and the sum extends over all square divisors $\mu\mu$ of the number $m$. This immediately gives
\[
f(m)=f(p^\pi p'^{\pi'}p''^{\pi''}\ldots)=f(p^\pi)f(p'^{\pi'})f(p''^{\pi''})\ldots
\]
and
\[
f(p^\pi)=\psi(p^\pi)+\psi(p^{\pi-2})+\psi(p^{\pi-4})+\ldots
\]
the series being continued as long as the exponents $\pi,\ \pi-2,\ \pi-4\ldots$ do not become negative. When $p$ does not divide $2D$, this gives
\[
f(p^\pi)=1+\left(\frac{-D}p\right)+\left(\frac{-D}{pp}\right)+\ldots+\left(\frac{-D}{p^\pi}\right)
\]
and, in general, when $m$ is relatively prime to $2D$,
\[
f(m)=\Sigma\left(\frac{-D}n\right)
\]
where the sum extends over all divisors $n$ of the number $m$.

\Needspace{18\baselineskip}
\SourceLeaf{304}
These observations immediately establish the statements in the text (III, 1, 2, 3) about the number $(m)$, provided that the first is supplemented by the condition that $D$ must not be divisible by $pp$ (the determination of the number of classes was already reduced in Disqq. Arithm., art. 256, to the case in which $D$ is divisible by no square). At the same time, even without this restriction, the infinite series
\[
1+\frac{f(p)}p+\frac{f(pp)}{pp}+\frac{f(p^3)}{p^3}+\ldots
\]
is found to have the value
\[
\frac1{1-\frac1p}\quad\text{or}\quad\frac1{1-.\frac1p}\pdot\frac1{1-\left(\frac{-D}p\right)\frac1p}
\]
according as $2D$ is or is not divisible by the prime $p$.
\BookRule{.25\textwidth}
